\documentclass[aspectratio=169,11pt]{beamer}
% Shared preamble for the CogLean decks.
% Build the slides alone, or with speaker notes (see the Makefile):
%   pdflatex deck.tex                                   -> slides
%   pdflatex -jobname=deck-notes "\def\withnotes{}\input{deck.tex}"  -> slides + notes
\usetheme{default}
\usecolortheme{seahorse}
\setbeamertemplate{navigation symbols}{}
\setbeamertemplate{footline}[frame number]
\setbeamerfont{note page}{size=\small}
\ifdefined\withnotes
  \setbeameroption{show notes}
\fi
\usepackage[T1]{fontenc}
\usepackage{lmodern}
\usepackage{booktabs}
\usepackage{amsmath,amssymb}
\usepackage{tikz}
\usetikzlibrary{positioning,arrows.meta,shapes.misc}
\tikzset{
  box/.style={draw, rounded corners, align=center, minimum height=8mm, inner sep=4pt, fill=blue!6},
  arr/.style={-{Stealth[length=2mm]}, thick}
}
\newcommand{\code}[1]{\texttt{#1}}
\newcommand{\has}{\textbf{has}}
\newcommand{\lacks}{\textcolor{red!70!black}{lacks}}
\newcommand{\never}{\textcolor{red!70!black}{\textbf{never}}}
\newcommand{\canhold}{can hold}

% Lean source in slides: fvextra Verbatim, with the Unicode Lean uses mapped
% to math symbols (pdflatex).
\usepackage{fvextra}
\fvset{fontsize=\scriptsize, frame=single, framesep=2mm, rulecolor=\color{gray!50},
  breaklines=true, breakanywhere=false}
\DeclareUnicodeCharacter{2192}{\ensuremath{\to}}
\DeclareUnicodeCharacter{2190}{\ensuremath{\leftarrow}}
\DeclareUnicodeCharacter{2194}{\ensuremath{\leftrightarrow}}
\DeclareUnicodeCharacter{2200}{\ensuremath{\forall}}
\DeclareUnicodeCharacter{2203}{\ensuremath{\exists}}
\DeclareUnicodeCharacter{27E8}{\ensuremath{\langle}}
\DeclareUnicodeCharacter{27E9}{\ensuremath{\rangle}}
\DeclareUnicodeCharacter{2260}{\ensuremath{\neq}}
\DeclareUnicodeCharacter{2227}{\ensuremath{\wedge}}
\DeclareUnicodeCharacter{2228}{\ensuremath{\vee}}
\DeclareUnicodeCharacter{00AC}{\ensuremath{\neg}}
\DeclareUnicodeCharacter{2264}{\ensuremath{\le}}
\DeclareUnicodeCharacter{2208}{\ensuremath{\in}}
\DeclareUnicodeCharacter{2209}{\ensuremath{\notin}}
\DeclareUnicodeCharacter{00D7}{\ensuremath{\times}}
\DeclareUnicodeCharacter{03B1}{\ensuremath{\alpha}}
\DeclareUnicodeCharacter{03B2}{\ensuremath{\beta}}
\DeclareUnicodeCharacter{03BC}{\ensuremath{\mu}}
\DeclareUnicodeCharacter{03B8}{\ensuremath{\theta}}
\DeclareUnicodeCharacter{2218}{\ensuremath{\circ}}
\DeclareUnicodeCharacter{25B8}{\ensuremath{\blacktriangleright}}
\DeclareUnicodeCharacter{00B7}{\ensuremath{\cdot}}
\DeclareUnicodeCharacter{2026}{\ldots}


% Generated by gen_story.py from the other five decks; do not edit. The core
% slide of each, copied verbatim, in an order that tells the story in one talk.

\title{CogLean in five slides}
\subtitle{Why prove, what is proved, how, is it faithful, and where it goes}
\author{Jeff Shrager}
\date{October 2026}

\begin{document}

\begin{frame}
  \titlepage
\end{frame}
\note{
\begin{itemize}
\item One straight-through talk: the core slide from each of the five CogLean decks. Each slide's notes begin with the transition, then the original notes.
\end{itemize}
}

\begin{frame}{Three kinds of evidence a theory can give}
  \small
  \begin{tabular}{p{0.3\linewidth}p{0.3\linewidth}p{0.3\linewidth}}
    \toprule
    \textbf{One run} & \textbf{Every run} & \textbf{Every model satisfying the postulates}\\
    \midrule
    one model, one parameter setting & one model, all its executions & all models the theory allows\\
    the theory \emph{can} do it & that model \emph{always} does it & the postulates \emph{entail} it\\
    simulation & model checking, sweeps & proof\\
    \bottomrule
  \end{tabular}
  \bigskip

  CogLean is built around the third kind, and keeps the other two in the same language.
\end{frame}
\note{
\textbf{Transition.} Why prove things about cognitive theories at all.
\begin{itemize}
\item A theorem quantified over every model satisfying the postulates says which assumptions matter (the proof uses some and not others), what is robust, and whether a claim is possibility or necessity.
\item Simulation is still essential: it explores, suggests conjectures, and fits data. Proof settles which conjectures follow from which assumptions. Experiments with people decide whether the assumptions are right.
\end{itemize}
}

\begin{frame}{Confident errors}
  Let a familiar problem $q$ have \emph{sharp} associations (only its correct sum clears the criterion), and let $x$, with a different sum, be recognized as $q$. Then:
  \begin{enumerate}
    \item on $q$, whole-problem retrieval never errs;
    \item on $x$, every whole-problem answer is wrong;
    \item $x$ uses the backup exactly as rarely as $q$ does: $(1-R_q)^n$.
  \end{enumerate}
  \bigskip
  In DOA alone, errors and backup use rise together. Confusion breaks the tie: a retrieved, wrong answer with no extra backup use.
\end{frame}
\note{
\textbf{Transition.} What gets proved: a prediction that neither EPAM nor the distributions-of-associations model makes alone, holding for every parameter setting. (Full talk: recognition-retrieval.)
\begin{itemize}
\item This is the prediction neither theory makes alone.
\item Two cautions. The model has no timing, so ``fast'' would need a timing model. And the errors are frequent only if retrieval on $q$ usually succeeds ($R_q$ near 1); if $R_q = 0$, every trial uses the backup.
\end{itemize}
}

\begin{frame}[fragile]{Confident errors: the proof is a composition}
\begin{Verbatim}
theorem confident_errors (correct : P → A) (x q : P) (hT : ∀ y, total as (str y) ≠ 0)
    (hbx : mass (backup x) = 1) (hbq : mass (backup q) = 1)
    (hq : retrieve (code q) = pure (some (code q)))
    (hx : retrieve (code x) = pure (some (code q)))
    (hs : Sharp str cc (code q) (correct q)) (hne : correct x ≠ correct q) :
    prob (recognizeRetrieve code retrieve str as cc n backup q) (wholeWrong correct q) = 0 ∧
      prob (recognizeRetrieve code retrieve str as cc n backup x) (wholeWrong correct x) =
        prob (recognizeRetrieve code retrieve str as cc n backup x) isWhole ∧
      prob (recognizeRetrieve code retrieve str as cc n backup x) isBackup =
        prob (recognizeRetrieve code retrieve str as cc n backup q) isBackup := by
  refine ⟨confusion_gives_neighbours_answer code retrieve str as cc n backup correct q q (hT _) hq hs,
    wrong_eq_whole_of_sharp code retrieve str as cc n backup correct x _ _ (hT _) hx hs hne, ?_⟩
  rw [backup_prob_certain code retrieve str as cc n backup x _ hx hbx hT,
    backup_prob_certain code retrieve str as cc n backup q _ hq hbq hT]
\end{Verbatim}
\end{frame}
\note{
\textbf{Transition.} How it is proved: the actual Lean proof is five lines, because the work is factored into general lemmas below it. (Full talk: recognition-retrieval-lean.)
\begin{itemize}
\item Three conjuncts, three earlier lemmas. The first is the neighbour's-answer lemma applied to $q$ itself; the second says that on sharp associations, every stated answer to $x$ is wrong; the third rewrites both backup probabilities to the same closed form $(1-R_q)^n$, after which the goal is an identity.
\item Note the hypotheses: retrieval is certain (\code{= pure (some \ldots)}). That is exactly what ACT-R never satisfies, and what the laws later supply.
\end{itemize}
}

\begin{frame}{Results: no failures}
  \small
  \begin{center}
  \resizebox{\linewidth}{!}{%
  \begin{tabular}{lllrr}
    \toprule
    Check & Lean theorem & kind & tests & failures\\
    \midrule
    Lean exact & \code{\#eval MD.trial} & prob & 243 & 0\\
    P1 & \code{columns\_prob} & prob & 394 & 0\\
    P2 & \code{whole\_depends\_only\_on\_route} & prob & 8543 & 0\\
    P3 & \code{confusion\_gives\_neighbours\_answer} & exact & 200 & 0\\
    P4 & \code{confident\_errors} & exact+prob & 222 & 0\\
    P7 & \code{learning\_preserves\_trial} & exact & 74692 & 0\\
    P9 & \code{columns\_place\_value} & exact & 248800 & 0\\
    P11 & \code{discrimination\_removes\_confusion} & exact & 91090 & 0\\
    P6a, P6b, P14, P17, P18 & & & 0 & \emph{vacuous}\\
    \bottomrule
  \end{tabular}}
  \end{center}
  \medskip
  \normalsize
  The prose and the formalization agree on recognition, EPAM learning, retrieval, both backups, and carrying.
\end{frame}
\note{
\textbf{Transition.} Is the formalization faithful? A Python simulator written blind from a prose specification, against predictions preregistered before it existed. (Full talk: blind-simulation.)
\begin{itemize}
\item A selection of rows; the full table (22 checks) is in \code{MultiDigit/sim/RESULTS.md}.
\item ``Lean exact'': 243 outcome frequencies on the four shared instances, each within tolerance of Lean's exact probability.
\end{itemize}
}

\begin{frame}{The matrix}
  \small
  \begin{tabular}{lcccccc}
    \toprule
    & Certainty & Stability & Order-free & Recency & Stored-only & Determinate\\
    \midrule
    Discrimination & \has & \has & \lacks & \lacks & \has & \has\\
    Activation & \lacks & \lacks & \lacks & \has & \has & \lacks\\
    Strength & \lacks & \lacks & \has & \lacks & \has & \lacks\\
    Global matching & \lacks & \lacks & \has & \lacks & \lacks & \has\\
    \bottomrule
  \end{tabular}
  \bigskip
  \begin{itemize}\small
    \item \has: a general proof, for every parameter choice meeting stated conditions.
    \item \lacks: a counterexample, computed exactly on the eight 3-bit codes.
    \item Every law is had by some family and lacked by another; every family has its own row. A script checks both.
  \end{itemize}
\end{frame}
\note{
\textbf{Transition.} Where it goes: instead of formalizing models one by one, laws of memory, and a matrix of which families have which law, every cell a theorem. Experiments can then be derived from laws that disagree. (Full talk: laws-matrix-experiments.)
\begin{itemize}
\item Examples of ``has'': EPAM's certainty needs only that the items be noticeably different; ACT-R's recency holds for any positive mismatch discount and failure weight (the later item is retrieved exactly twice as often); strength memory and MINERVA are order-free for every parameter choice.
\item A ``lacks'' cell says only that the law fails somewhere. Mapping the parameter regions where a family does have it is open.
\end{itemize}
}

\end{document}
